% ============================================================
% Pose-Differentiable Body Twins (RIHB-in-the-loop)
% First-principles companion to rib.tex and q_complement.tex.
% Intended home for the deployable parts: a JSAC2 SI submission
% on Digital Twins for Wireless Networks (Q4 2026 issue).
% Author: R. Wydaeghe
% ============================================================

\documentclass[11pt,a4paper]{article}

\usepackage[margin=2.4cm]{geometry}
\usepackage[utf8]{inputenc}
\usepackage[T1]{fontenc}
\usepackage{lmodern}
\usepackage{microtype}
\linespread{1.04}

\usepackage{amsmath,amssymb,amsthm,mathtools}
\usepackage{bm}
\usepackage{mathrsfs}
\usepackage{bbm}
\usepackage{cancel}
\usepackage{xfrac}
\allowdisplaybreaks[2]

\usepackage{empheq}
\usepackage[most]{tcolorbox}
\tcbset{highlight math style={%
    enhanced, colframe=black!70, colback=gray!6,
    arc=2pt, boxrule=0.5pt, left=3pt, right=3pt}}

\usepackage{siunitx}
\sisetup{detect-all, range-phrase = {--}, range-units = single,
         per-mode = symbol, inter-unit-product = {\,}}

\usepackage{enumitem}
\usepackage{xcolor}
\usepackage{xspace}
\usepackage{fancyhdr}
\pagestyle{fancy}
\fancyhf{}
\cfoot{\thepage}
\renewcommand{\headrulewidth}{0pt}

\usepackage{titlesec}
\titleformat{\section}{\large\bfseries}{\thesection}{1em}{}
\titleformat{\subsection}{\normalsize\bfseries}{\thesubsection}{1em}{}

\usepackage[colorlinks=true, linkcolor={blue!55!black},
            citecolor={green!45!black}, urlcolor={blue!65!black}]{hyperref}
\usepackage[nameinlink,capitalize]{cleveref}

\usepackage{aliascnt}
\theoremstyle{definition}
\newtheorem{theorem}{Theorem}[section]
\newaliascnt{proposition}{theorem}
\newtheorem{proposition}[proposition]{Proposition}
\aliascntresetthe{proposition}
\newaliascnt{lemma}{theorem}
\newtheorem{lemma}[lemma]{Lemma}
\aliascntresetthe{lemma}
\newaliascnt{corollary}{theorem}
\newtheorem{corollary}[corollary]{Corollary}
\aliascntresetthe{corollary}
\newaliascnt{definition}{theorem}
\newtheorem{definition}[definition]{Definition}
\aliascntresetthe{definition}
\newaliascnt{remark}{theorem}
\newtheorem{remark}[remark]{Remark}
\aliascntresetthe{remark}
\crefname{proposition}{Proposition}{Propositions}
\Crefname{proposition}{Proposition}{Propositions}
\crefname{lemma}{Lemma}{Lemmas}\Crefname{lemma}{Lemma}{Lemmas}
\crefname{corollary}{Corollary}{Corollaries}\Crefname{corollary}{Corollary}{Corollaries}
\crefname{definition}{Definition}{Definitions}\Crefname{definition}{Definition}{Definitions}
\crefname{remark}{Remark}{Remarks}\Crefname{remark}{Remark}{Remarks}

% Macros aligned with monograph, q_complement, rib.tex
\newcommand{\khat}{\hat{\bm{k}}}
\newcommand{\nhat}{\hat{\bm{n}}}
\newcommand{\rhat}{\hat{\bm{r}}}
\newcommand{\rr}{\mathbf{r}}
\newcommand{\EE}{\mathbf{E}}
\newcommand{\Sinc}{S_{\mathrm{inc}}}
\newcommand{\Sab}{S_{\mathrm{ab}}}
\newcommand{\Sref}{S_{\mathrm{re}}}
\newcommand{\Sinn}{S_{\mathrm{in}}}
\newcommand{\Smis}{S_{\mathrm{mi}}}
\DeclareMathOperator{\ReLU}{ReLU}
\DeclareMathOperator{\RE}{Re}
\DeclareMathOperator{\IM}{Im}
\DeclareMathOperator{\tr}{tr}
\DeclareMathOperator*{\argmin}{arg\,min}
\DeclareMathOperator*{\argmax}{arg\,max}
\newcommand{\diff}{\mathrm{d}}
\newcommand{\ident}{\mathbbm{1}}
\newcommand{\Reals}{\mathbb{R}}
\newcommand{\Complex}{\mathbb{C}}
\newcommand{\half}{\tfrac{1}{2}}
\newcommand{\Gtil}{\tilde{\mathbf{G}}}
\newcommand{\gtil}{\tilde{\mathbf{g}}}
\newcommand{\Fn}{\mathbf{F}_{\!n}}
\newcommand{\Rn}{\mathbf{R}_{\!n}}
\newcommand{\Qabs}{\mathbf{Q}_{\mathrm{ab}}}
\newcommand{\Qref}{\mathbf{Q}_{\mathrm{re}}}
\newcommand{\Qinn}{\mathbf{Q}_{\mathrm{in}}}
\newcommand{\Qmis}{\mathbf{Q}_{\mathrm{mi}}}

% RIB / RIHB macros
\newcommand{\kin}{\hat{\bm{k}}_{\mathrm{i}}}
\newcommand{\kout}{\hat{\bm{k}}_{\mathrm{o}}}
\newcommand{\Frib}{\mathbf{F}_{\mathrm{b}}}
\newcommand{\sigb}{\sigma_{\mathrm{b}}}
\newcommand{\Qbist}{\mathbf{Q}_{\mathrm{bi}}}
\newcommand{\hh}{\mathbf{h}}
\newcommand{\Jmat}{\mathbf{J}}
\newcommand{\bvec}{\mathbf{b}}
\newcommand{\xvec}{\mathbf{x}}
\newcommand{\thetab}{\bm{\theta}}
\newcommand{\betab}{\bm{\beta}}
\newcommand{\Mmesh}{\mathcal{M}}
\newcommand{\LBS}{\mathrm{LBS}}

\title{\textbf{Pose-Differentiable Body Twins:\\
First Principles for User-Adaptive\\
mmWave Capacity and Exposure}}
\author{Robin Wydaeghe\\
 Ghent University \& imec\\
 \texttt{robin.wydaeghe@ugent.be}}
\date{\today}

\makeatletter
\renewcommand{\maketitle}{%
  \thispagestyle{plain}%
  \begin{center}
    \rule{\linewidth}{1pt}\\[0.4em]
    {\LARGE\bfseries \@title}\\[0.5em]
    {\large \@author}\\[0.3em]
    {\small \@date}\\[0.4em]\rule{\linewidth}{1pt}
  \end{center}
  \vspace{0.5em}
}
\makeatother

\begin{document}
\maketitle

\begin{abstract}
A user with a smartphone is, at mmWave frequencies, an electromagnetic
object that the network can both \emph{model} and \emph{negotiate
with}.  This note develops the simplest non-trivial setup that makes
both halves of that statement quantitative.  The body is a closed
parametric surface $\Sigma(\thetab,\betab)$ with pose
$\thetab\in\Reals^{72}$ given by an SMPL-X estimator running on the
device.  The exposure operator $\Qabs(\thetab)$ and the bistatic
scattering matrix $\Frib(\thetab,\kout)$ inherit the same
factorisation as the body twin in the JSAC paper, and are
\emph{differentiable in $\thetab$} almost everywhere.  This makes the
user a low-dimensional reconfigurable surface in the propagation
graph, with one knob (gait) that the user controls and the array
predicts.  Closing the loop, the array sends the user a per-slot
gradient suggestion $\nabla_\thetab\,\mathcal{L}$ and the user, at a
\SI{1}{Hz} cadence and a comfort budget, chooses to either follow it
or not.  When followed, the body shapes its own channel; when not,
the array continues to use $\Qabs(\thetab)$ to keep the user inside
the basic restriction.  We formalise this as a min-max problem on a
joint capacity-exposure objective, refresh the per-direction
pseudo-Brewster Pareto bound of \texttt{rib.tex} into its
single-user user-adaptive form, and lay out the information-theoretic
budget of the closed loop ($\sim$~\SI{18}{kbps} bidirectional, four
orders of magnitude below the data rate it negotiates over).  The
hero claim is the empirical one and is left to a companion experiment
document: in a body-shadowed line-of-sight scenario, sub-degree pose
adjustments suffice to traverse a Pareto frontier whose dynamic range
on capacity is set by the body's blockage cross-section, and whose
dynamic range on dose is set by $(1-T_0)/T_0$, the body-universal
mmWave reflection-to-absorption ratio.
\end{abstract}

\vspace{0.3em}
\noindent\textit{Style note.}~Like its siblings
\texttt{rib.tex} and \texttt{q\_complement.tex}, this is a wandering
companion, not a theorem-proof paper.  Theorems are theorems where
the proof is self-contained or already in a sibling document;
propositions are claims I want a reviewer to take or leave on
sketches; conjectures are flagged.  The intended home for the
deployable parts is \S\,IV of a JSAC2 SI submission.  The intended
home for everything else is the next companion.

\tableofcontents
\bigskip

% ============================================================
\section{Why a fifth body-twin paper}\label{sec:why}
% ============================================================

The body-twin papers in the lineage \texttt{(paper, paper\_v2,
paper\_jsac\_v1\dots v5)} converged on a strong and, in retrospect,
narrow result: under the operational \mbox{5G NR FR2} per-stream
MCS27 cap, regularised zero-forcing on a tens-of-element panel
saturates the rate cap with $\sim$\SI{12}{dB} of SINR margin, and a
one-line scaling on that precoder enforces per-body exposure
compliance by construction.  The body twin's role inside that result
is to compute one scalar
$\alpha=\sqrt{\eta_S\,\min_u L_{\mathrm{RL}}^{(u)}/P_{\mathrm{ab}}^{(u)}}$
per slot.  The twin is structurally a per-body absorbed-power
estimator, and the deployment value of the twin over a worst-case
Cauchy envelope is real but small.

The JSAC SI on Digital Twins for Wireless Networks asks for
something different.  Its scope keywords are
\emph{closed-loop optimization}, \emph{differentiable control},
\emph{twin-in-the-loop architectures}, and
\emph{cross-layer feedback}.  None of these describe a one-line
amplitude-scaling of regularised ZF.  This note re-targets the body
twin from \emph{constraint provider} to \emph{control variable}: the
user's pose $\thetab$ becomes the channel-shaping degree of freedom
the network negotiates with, and the body twin is the model that
makes the negotiation tractable.

There are three things this re-targeting buys, and one thing it gives
up.  The three things:

\begin{enumerate}[leftmargin=1.6em,itemsep=3pt]
\item \textbf{The twin becomes load-bearing.}  Without
$\thetab$-differentiability, no closed loop is possible.  Pose, scene,
and exposure must all sit inside the same gradient.  This rules out
the worst-case envelope that the JSAC v5 result tacitly competes
against: a Cauchy bound ignores $\thetab$ entirely.
\item \textbf{The application is the user.}  ``Application-aware''
in the SI scope is a phrase that points naturally to robotics, XR,
and autonomy.  The user holding a smartphone is a flat-corner case
the literature has not picked up: the network is serving a session
\emph{to} the user's device while the user's body is in the
propagation graph.  Re-posing the user is the cheapest, lowest-power
reconfigurable surface action available; an upright human standing
still needs to expend $\sim$\SI{1}{W} to maintain posture, against
which a few-degree torso turn is free.
\item \textbf{The exposure-vs-capacity Pareto becomes a UI.}
The companion document \texttt{rib.tex} (\S\,8) lifted in
\cref{sec:pareto} below shows that body-mediated capacity gain is
upper-bounded by
absorbed dose with a body-universal slope of $T_0/(1-T_0)\approx
0.32$ at \SI{28}{GHz}.  This converts ``show me my SAR while I
download'' from a vanity meter into the price tag of the capacity
the user just bought.  The user can read the gradient, decide
whether to take it, and the receipt is honest by physics.
\end{enumerate}

The thing it gives up is the deployment-without-consent path.  The
v5 paper is a story the operator can run with no UE-side cooperation
at all.  This note's spine assumes the user opted in.  Whether the
opt-in fraction in deployment is 1\% or 60\% is a UX question.  We
will argue that the architecture survives at small opt-in fractions
because every co-operating user lifts \emph{others'} channels too via
the body-as-passive-RIS effect of \texttt{rib.tex} \S\,9.

% ============================================================
\section{Setup and notation}\label{sec:setup}
% ============================================================

The setup is the JSAC paper's, plus pose.

\subsection{Geometry}

A single base station hosts an array of $M$ elements at known
positions $\{\rr_j\}_{j=1}^M$.  In the running examples we take a
$M=64$ uniform rectangular array (URA) with half-wavelength spacing
at carrier frequency $f_c=\SI{28}{GHz}$
($\lambda\approx\SI{1.07}{cm}$).  The array radiates the precoder
$\xvec\in\Complex^M$ with sum-power constraint
$\|\xvec\|^2\le P_{\mathrm{tx}}$.

A single user is described by a parametric body model.  We use
SMPL-X (Pavlakos et al.\ 2019) with shape vector $\betab\in\Reals^{10}$
(identity, fixed once per user) and pose vector
$\thetab\in\Reals^{72}$ (root orientation plus 21 axis-angle joint
rotations plus a few facial/hand DOF that we will ignore).  The body
mesh is the closed orientable surface
\begin{equation}\label{eq:mesh}
  \Sigma(\thetab,\betab)
   =\Mmesh\bigl(\LBS(\thetab,\betab,\mathbf{T})\bigr),
\end{equation}
where $\mathbf{T}\in\Reals^{V\times 3}$ is the SMPL-X template with
$V=\num{10475}$ vertices and $\Mmesh$ assembles the $F=\num{20908}$
triangles of the topology.  $\LBS$ is the linear blend skinning
operator: each vertex's deformed position is a
weight-blended composition of joint rotations and translations.  The
body's outward normal at a triangle centroid is
$\nhat(\rr;\thetab,\betab)$.

The user holds a smartphone receiver at a position $\rr_p$ that is
itself a function of pose: when the user moves their arm, $\rr_p$
moves.  We model the phone as a single-element receiver whose pose
relative to the wrist is fixed; the wrist position is one of the
SMPL-X joints.  Hence $\rr_p=\rr_p(\thetab,\betab)$ is also smooth
in $\thetab$ almost everywhere.

\subsection{Path-space factorisation}

The propagation between the array and any spatial probe (the body's
surface, the phone's location) is described by $N$ multipath
contributions $\{j(n),\khat_n,\bm\psi_n,\tau_n\}_{n=1}^N$ as in the
JSAC paper, with $j(n)\in\{1,\dots,M\}$ the source element, $\khat_n$
the unit propagation direction at the probe, $\bm\psi_n\in\Complex^3$
the polarisation amplitude, and $\tau_n$ the path delay.  Path
construction is delegated to a ray tracer (Sionna or DiffeRT).  At
the body, the per-element steering response across the lit area
factors as $\mathbf{J}^T\mathbf{a}(\khat_n)$, where $\mathbf{J}\in
\Complex^{M\times M}$ encodes element-to-element phase relationships
and $\mathbf{a}(\khat)$ is the array manifold; the JSAC paper denotes
the assembled per-body channel
$\Gtil^{(u)}(\rr;\thetab,\betab)\in\Complex^{3\times M}$.

The two operators that recur throughout this note:
\begin{align}
  \Qabs(\thetab,\betab)
  &= \int_{\Sigma(\thetab,\betab)}
     \Gtil(\rr;\thetab,\betab)^H\,\Gtil(\rr;\thetab,\betab)\,
     \mu(\rr;\thetab)\,\diff A(\rr),
  \label{eq:Qabs-def}\\[0.2em]
  \Qbist(\thetab,\betab,\kout)
  &= \frac{1}{Z_0}\,\Frib(\thetab,\betab,\kout)^H\,
     \Frib(\thetab,\betab,\kout),
  \label{eq:Qbist-def}
\end{align}
with $\mu(\rr;\thetab)=\nhat(\rr;\thetab)\cdot(-\khat)$ the
incidence cosine, $Z_0$ the free-space impedance, and
$\Frib(\thetab,\betab,\kout)$ the body scattering matrix from
\texttt{rib.tex} (\S\,4, Definition~``body scattering matrix'').  The
dependence on $\kout$ in $\Qbist$ is the per-direction
angular-density of the reflection operator $\Qref$, and integrating
$\Qbist$ over the back-hemisphere recovers $\Qref$ to within the
$5\%$ Kirchhoff residual (\texttt{rib.tex} \S\,5,
``bookkeeping reconciliation'').

\subsection{What is new in this note}

Both operators were treated as static functions of the body's pose
in the JSAC v1--v5 papers: pose was an estimated parameter, the twin
re-evaluated $\Qabs$ at the new estimate every \SI{100}{ms}, and the
precoder side used the result.  In this note, $\thetab$ becomes a
\emph{decision variable}.  Every place pose was an input becomes a
place pose is a knob.  The mathematics is the same; the question
posed of it is different.

% ============================================================
\section{The user as a reconfigurable surface}\label{sec:user-recon}
% ============================================================

The conceptual move is to treat the user the way RIS literature
treats a reconfigurable surface, with the modifications that physics
imposes on the analogy.  This subsection makes those modifications
explicit so we do not oversell.

\subsection{The body-state vector}\label{sec:state}

\begin{definition}[User control vector]\label{def:phi}
The user control vector
$\bm\varphi=(\thetab,\rr_{\mathrm{root}},\bm\omega_{\mathrm{root}})
\in\Reals^{78}$
collects the SMPL-X pose, the root translation in world coordinates,
and the root orientation.  We will use $\thetab$ as shorthand for
$\bm\varphi$ when the root pose is fixed (e.g.\ a seated user at a
desk); the more general theory uses $\bm\varphi$.
\end{definition}

Not all $78$ DOF are negotiable by the user at the same cost.  The
following is folklore among biomechanics and pose-tracking
practitioners; we list it because it shows up in the loss design.

\begin{enumerate}[leftmargin=1.6em,itemsep=2pt]
\item \emph{Free DOF.}  Eyebrow lift, finger curl, jaw position.  No
metabolic cost, but they move negligible mesh area and have no
appreciable mmWave effect.
\item \emph{Comfortable DOF.}  Head yaw/pitch, torso yaw, arm and
forearm joint angles, hand grip orientation.  Sub-degree adjustments
are imperceptible; few-degree adjustments are easy.  These are the
DOF that move the body's projected area in the direction of the
LOS path significantly.
\item \emph{Costly DOF.}  Spine bend, kneeling, leaning beyond the
support polygon.  Metabolically expensive, socially conspicuous,
and the user will ignore suggestions in this category.
\item \emph{Rigid DOF.}  Root translation in world coordinates
(walking).  Available but slow and negotiated externally to the
session.
\end{enumerate}

We model the cost of the move from baseline pose $\thetab^{(0)}$ to
proposed pose $\thetab$ as a separable convex function
\begin{equation}\label{eq:movement-cost}
  C(\thetab-\thetab^{(0)})
  = \sum_{i=1}^{72}c_i\,(\theta_i-\theta_i^{(0)})^2,
\end{equation}
with $c_i$ small for comfortable DOF, large for costly DOF, infinite
for joint-limit violations.  In the experiments of the companion
document the comfortable-DOF coefficient is calibrated such that
$C\le 1$ corresponds to a \SI{5}{\degree} torso rotation.

\subsection{Update rate and actuation manifold}

Pose updates are pulled to the user.  Two cadences matter.

\paragraph{Internal pose drift.}  Even a still user has pose noise:
breathing, postural sway, micro-arm-motion.  Across \SI{100}{ms}
windows, $\mathrm{std}(\thetab)\sim$ a few tenths of a degree per
joint at the torso, more at the wrist.  This is the cadence the
\texttt{virtual\_imu.py} pipeline tracks at and is what the JSAC v5
paper budgets the twin against.

\paragraph{Negotiated pose move.}  When the BS suggests a move, the
user reads it (visual notification on the phone) and, if they take
it, the move completes in \SI{200}{ms} to \SI{2}{s}.  The control
loop runs at \SI{1}{Hz}.  This is the cadence we design the joint
optimiser against in \cref{sec:loss}; the inner pose-drift cadence
is handled by the JSAC v5 machinery unchanged.

The actuation manifold is the user's natural pose distribution
truncated to comfortable DOF.  We will model it implicitly via the
weights $c_i$ in \eqref{eq:movement-cost} rather than as an explicit
manifold; this is enough for the gradient updates to respect joint
limits with a barrier term $\sum_i\beta_i\log(1-(\theta_i/\theta_i^{\max})^2)$
absorbed into $C$.

% ============================================================
\section{Differentiability of the body channel}\label{sec:diff}
% ============================================================

This section is the load-bearing analytical claim of the note.  Every
sub-claim has either a one-line proof or a reference to a sibling
document where the technical work is done.

\subsection{Linear blend skinning is smooth a.e.\ in pose}

\begin{lemma}[LBS smoothness]\label{lem:lbs-smooth}
The map $\thetab\mapsto\Sigma(\thetab,\betab)$ defined by
\eqref{eq:mesh} is a piecewise-smooth function of $\thetab$ on
$\Reals^{72}$, with derivative discontinuities only at pose
configurations where two distinct joints' weight blends are exchanged
by a continuous deformation (a measure-zero set).  Equivalently,
each vertex position $\rr_v(\thetab,\betab)$ is given by
\begin{equation}\label{eq:lbs}
  \rr_v(\thetab,\betab)
  =\sum_{j=1}^{J}w_{vj}\bigl(\mathbf{R}_j(\thetab)\,
                              (\mathbf{T}_v+\mathbf{B}_v(\betab)
                                              -\mathbf{j}_j)
                              +\mathbf{j}_j+\mathbf{t}_j(\thetab)\bigr),
\end{equation}
where $w_{vj}$ are the SMPL-X skinning weights, $\mathbf{R}_j(\thetab)
\in SO(3)$ is the world rotation of joint $j$ (a polynomial in
$\sin\theta_{j,\cdot},\cos\theta_{j,\cdot}$), $\mathbf{j}_j$ is the
canonical joint position, and $\mathbf{t}_j(\thetab)$ is the joint's
forward-kinematic translation.  Both $\mathbf{R}_j$ and
$\mathbf{t}_j$ are $C^\infty$ in $\thetab$; the blend
$\sum_jw_{vj}(\cdot)$ preserves smoothness.
\end{lemma}

\noindent The proof is in any LBS reference (the SMPL paper, or
Lewis-Cordner-Fong); we re-state it because the
$C^\infty$-everywhere claim is what we need below, and the
exchange-of-joints non-smoothness has zero measure on natural pose
distributions.  Triangle normals
$\nhat_t(\thetab,\betab)$ are the cross product of two vertex
differences, hence $C^\infty$ in $\thetab$ wherever the triangle is
non-degenerate.

\subsection{The triangle integrand}

For one path $n$ with direction $\khat_n$ and polarisation
$\bm\psi_n$, the contribution to $\Qabs$ from triangle $t$ at
centroid $\rr_t$ with area $A_t$ and normal $\nhat_t$ is
\begin{equation}\label{eq:int-tri}
  q_{t,n}(\thetab) = A_t(\thetab)\,\mu_{t,n}(\thetab)\,
                     \mathbf{F}_n(\rr_t;\nhat_t,\khat_n)^H
                     \mathbf{F}_n(\rr_t;\nhat_t,\khat_n)\,H(\mu_{t,n}),
\end{equation}
with $\mathbf{F}_n$ the per-path Fresnel transmission operator from
\texttt{q\_complement.tex} \S\,3, and $H(\cdot)$ the Heaviside
visibility gate enforcing front-facing.  The Heaviside is the only
non-smooth ingredient.

\begin{lemma}[Triangle smoothness]\label{lem:tri-smooth}
Away from $\mu_{t,n}=0$, the map
$\thetab\mapsto q_{t,n}(\thetab)$ is $C^\infty$ on the LBS
smooth set of \cref{lem:lbs-smooth}.
\end{lemma}

\noindent This is immediate: $A_t=\half\|(\rr_2-\rr_1)\times
(\rr_3-\rr_1)\|$ is polynomial in vertex positions; $\mu_{t,n}
=\nhat_t\cdot(-\khat_n)$ is bilinear in normal and direction; the
Fresnel operator $\mathbf{F}_n$ is rational in $\mu$ (and smooth as a
function of pose at fixed tissue); compositions of $C^\infty$ maps
are $C^\infty$.

\subsection{Visibility smoothing}\label{sec:gelu}

The Heaviside discontinuity at $\mu_{t,n}=0$ is the standard
silhouette artefact in differentiable rendering.  We use the same
GELU smoothing as TAP \S\,2.6:
\begin{equation}\label{eq:gelu}
  H_\epsilon(\mu) \equiv \tfrac{1}{2}\bigl(1+\mathrm{erf}(\mu/(\epsilon\sqrt{2}))\bigr),
\end{equation}
with $\epsilon$ chosen at one degree of incidence-angle width
($\epsilon\approx 0.017$) so that the smoothing zone is narrower than
typical pose noise.  This replaces the hard $H(\mu_{t,n})$ in
\eqref{eq:int-tri} by $H_\epsilon(\mu_{t,n})$.

\begin{remark}\label{rem:gelu-bias}
GELU smoothing introduces a one-sided bias in $\Qabs$: triangles at
near-grazing incidence contribute partially.  At
$\epsilon=0.017$, the bias on $\tr\Qabs$ is $\le 0.4\%$ for
Thelonious at \SI{28}{GHz}; well below the JSAC paper's $\sim 5\%$
calibration floor.  We treat the smoothed $\Qabs$ as the
operationally correct one and the hard $\Qabs$ as a sanity check.
\end{remark}

\subsection{The end-to-end Jacobian}

Composing \cref{lem:lbs-smooth,lem:tri-smooth} and \eqref{eq:gelu}:

\begin{theorem}[Pose differentiability of $\Qabs$]\label{thm:Q-diff}
For fixed shape $\betab$, fixed path set
$\{\khat_n,\bm\psi_n\}_n$, and the GELU-smoothed visibility
$H_\epsilon$, the map
$\thetab\mapsto\Qabs(\thetab,\betab)\in\Complex^{M\times M}$
defined by \eqref{eq:Qabs-def} is $C^\infty$ on the LBS smooth set.
The Jacobian factorises as
\begin{empheq}[box=\tcbhighmath]{equation}\label{eq:Q-jac}
  \frac{\partial\Qabs}{\partial\theta_i}
  \;=\;\sum_{t=1}^{F}\sum_{n=1}^{N}
        \underbrace{\frac{\partial q_{t,n}}{\partial\rr_v}}_{\text{Fresnel-side}}
        \cdot
        \underbrace{\frac{\partial\rr_v}{\partial\theta_i}}_{\text{LBS-side, (3, V)}},
\end{empheq}
summed over the three vertices $v$ of triangle $t$.  The Fresnel-side
factor is local to one triangle and one path; the LBS-side factor is
global and is exactly the Jacobian SMPL-X provides via its PyTorch
autograd graph.
\end{theorem}

\noindent Proof.  Compose \cref{lem:lbs-smooth} for the LBS map with
\cref{lem:tri-smooth} for the integrand and the smoothed-Heaviside
for visibility.  $\Qabs$ is the sum of $F\cdot N$ per-triangle
per-path contributions, each $C^\infty$ in pose; the sum is
$C^\infty$.  The chain rule then gives \eqref{eq:Q-jac}; the
Fresnel-side and LBS-side factors are exactly the two terms emitted
by per-triangle backpropagation against vertex positions and per-vertex
backpropagation against pose.  $\square$

\Cref{thm:Q-diff} is the contract that lets us replace the JSAC v5
``re-evaluate $\Qabs$ at the IMU's pose estimate every
\SI{100}{ms}'' with ``compute $\partial\Qabs/\partial\thetab$ on the
device.''  The same theorem holds for the bistatic operator
$\Qbist(\thetab,\kout)$: the proof goes through verbatim with the
extra outgoing phase $e^{-ik_0\kout\cdot\rr}$, which is smooth in
$\rr$ and hence smooth in $\thetab$.

\subsection{Cost of the Jacobian}

The chain rule in \eqref{eq:Q-jac} costs one Fresnel-side evaluation
per (triangle, path) and one LBS-side evaluation per (vertex, pose
DOF).  In numbers: $F\cdot N\sim\num{20000}\cdot 50=10^6$
Fresnel-side terms (each a $3\times M$ complex matrix, so $\sim
M^2\cdot 10^6=4\cdot 10^9$ flops at $M=64$); $V\cdot 72\sim
\num{750000}$ LBS-side terms (each a $3\times 1$ real vector,
trivial).  The Fresnel side dominates and is identical in structure
to the JSAC paper's $\Qabs$ build, which already runs at
\SI{7.5}{ms} per body on a single GPU at SMPL-X mesh size.  A
forward-mode JVP of the same compute graph costs the same order;
reverse-mode VJP costs \SI{15}{ms} per pose direction.  The full
Jacobian is one VJP per output entry, but for the gradient updates
of \cref{sec:loss} we only need
$\partial\,\mathrm{tr}\Qabs/\partial\thetab\in\Reals^{72}$ which is
one VJP of one scalar against $\thetab$, total cost $\sim
\SI{15}{ms}$ per slot on a phone-class GPU.

% ============================================================
\section{The joint capacity-exposure loss}\label{sec:loss}
% ============================================================

We can now state the optimisation problem.

\subsection{Single-user achievable rate}\label{sec:rate}

Let $\hh(\thetab)\in\Complex^M$ be the array-to-phone channel
for the user, including any body-mediated paths via
$\Frib(\thetab,\kout)$ as in \texttt{rib.tex} (\S\,10, eq.\ for
$\bvec_{u\to\mathrm{UE}}$).
With single-stream MRT-MMSE precoding $\xvec=\sqrt{P_{\mathrm{tx}}}
\hh(\thetab)/\|\hh(\thetab)\|$, the SINR at the phone is
\begin{equation}\label{eq:sinr}
  \mathrm{SINR}(\thetab)
  =\frac{P_{\mathrm{tx}}\,\|\hh(\thetab)\|^2}{N_0\,B},
\end{equation}
and the achievable rate, capped by the operational MCS27 ceiling, is
\begin{equation}\label{eq:rate-mcs}
  R(\thetab)=B\cdot\min\!\bigl\{\log_2(1+\mathrm{SINR}(\thetab)),\;
  R_{\max}/B\bigr\},
  \qquad R_{\max}/B = \SI{7.4}{bps/Hz}.
\end{equation}
The $\min$ structure is the operational fact the JSAC v5 paper
foregrounds: above $\sim$\SI{22}{dB} SINR, more capacity is not
purchaseable per stream.

\subsection{Per-body exposure}

The user's per-stream absorbed power is
\begin{equation}\label{eq:Pabs-1}
  P_{\mathrm{ab}}(\thetab)=\xvec^H\Qabs(\thetab)\xvec,
\end{equation}
and the dose constraint is $P_{\mathrm{ab}}(\thetab)\le L_{\mathrm{RL}}$,
the basic restriction at the European reference level
$E_{\mathrm{RL}}=\SI{61}{V/m}$ (or \SI{3}{V/m} in the strict-Brussels
regime of JSAC v4--v5).

\subsection{The composite objective}

For $\lambda_E,\lambda_M>0$, define
\begin{empheq}[box=\tcbhighmath]{equation}\label{eq:loss}
  \mathcal{L}(\thetab,\xvec;\thetab^{(0)})
  = -R(\thetab,\xvec)
    \;+\;\lambda_E\,\bigl[P_{\mathrm{ab}}(\thetab,\xvec)
                            -L_{\mathrm{RL}}\bigr]_+
    \;+\;\lambda_M\,C(\thetab-\thetab^{(0)}).
\end{empheq}
The first term rewards capacity; the second term penalises exposure
overshoot; the third term penalises pose deviation from the user's
baseline.  All three are differentiable in $\thetab$ (the
$[\cdot]_+$ is GELU-smoothed at the same $\epsilon$ as the
visibility, for consistency).

\subsection{Inner-outer structure}

The natural solver alternates:
\begin{align*}
  \xvec^{(\ell+1)} &= \argmin_{\xvec:\,\|\xvec\|^2\le P_{\mathrm{tx}}}
                       \mathcal{L}(\thetab^{(\ell)},\xvec;\thetab^{(0)}),\\
  \thetab^{(\ell+1)}&= \argmin_{\thetab:\,C(\thetab-\thetab^{(0)})\le \rho}
                       \mathcal{L}(\thetab,\xvec^{(\ell+1)};\thetab^{(0)}).
\end{align*}
The inner ($\xvec$) problem is a QCQP per the JSAC v5 machinery.
Under the operational MCS-cap regime, it is the projected-ZF closed
form.  The outer ($\thetab$) problem is a $72$-dimensional smooth
optimisation with one nonconvex objective; gradient descent with
Armijo line search converges to a local stationary point at the
\SI{15}{ms} cost per gradient step from \cref{sec:diff}.

\subsection{Coupling: pose and precoder cannot be optimised separately}

The naive question ``why not optimise the pose first, then the
precoder?'' has a textbook answer: the channel
$\hh(\thetab)$ depends on pose, so the precoder that maximises
$\|\hh(\thetab)\|$ is itself $\thetab$-dependent.  The implicit
function theorem applied to the precoder optimum gives
\begin{equation}\label{eq:dchain}
  \frac{\diff R}{\diff\thetab}\bigg|_{\xvec^*(\thetab)}
  = \frac{\partial R}{\partial\thetab}
    + \frac{\partial R}{\partial\xvec}\frac{\diff\xvec^*}{\diff\thetab}
  = \frac{\partial R}{\partial\thetab},
\end{equation}
because the inner problem's KKT conditions zero out
$\partial R/\partial\xvec$ at $\xvec^*$.  Hence the outer gradient on
$\thetab$ at fixed $\xvec=\xvec^*(\thetab)$ is the partial
derivative; the alternating scheme is exact.  This is the standard
``envelope theorem'' / Danskin's theorem, useful here because it
spares us from differentiating through the QCQP solver.

% ============================================================
\section{The pose-conditional Pareto bound}\label{sec:pareto}
% ============================================================

The Pareto theorem of \texttt{rib.tex} (\S\,8) bounds body-mediated
capacity gain by absorbed dose with a body-universal slope.  Here
we tighten that bound for the single-user user-adaptive case.

\subsection{Single-user body-mediated capacity}

For one user with one body in the channel, the body-mediated
contribution to the SINR is the squared magnitude of the
body-scattered branch $\bvec_{u\to\mathrm{UE}}$ from \texttt{rib.tex}
\S\,10, which factors as
$\bvec_{u\to\mathrm{UE}}=\Frib(\thetab,\kout_{u\to\mathrm{UE}})\xvec/
[4\pi r_{u\to\mathrm{UE}}]$.  Hence the body-mediated contribution to
the receive power is
\begin{equation}\label{eq:Pbody}
  P_{\mathrm{body}}(\thetab,\xvec)
  =\frac{1}{(4\pi r_{u\to\mathrm{UE}})^2}\,
    \xvec^H\Qbist(\thetab,\kout_{u\to\mathrm{UE}})\xvec.
\end{equation}

\subsection{The bound}

\begin{theorem}[Single-user pose-conditional Pareto]\label{thm:pareto-1}
Under the assumptions of \texttt{rib.tex} \S\,3 (high-frequency
PO, far-field observer, single bounce, pseudo-Brewster), at any
fixed pose $\thetab$ and at the receiver direction
$\kout_{u\to\mathrm{UE}}=(\rr_p(\thetab)-\rr_b(\thetab))/r_{u\to\mathrm{UE}}$,
\begin{empheq}[box=\tcbhighmath]{equation}\label{eq:pareto-1}
  P_{\mathrm{body}}(\thetab,\xvec)
  \;\le\;
  \frac{1}{(4\pi r_{u\to\mathrm{UE}})^2}\,\frac{1-T_0}{T_0}\,
  \mathcal{V}_{\mathrm{out}}(\thetab)\,
  P_{\mathrm{ab}}(\thetab,\xvec)
  \;+\;O\!\bigl(5\%\bigr),
\end{empheq}
where $\mathcal{V}_{\mathrm{out}}(\thetab)\in[0,1]$ is the fraction
of the body's lit area visible from the phone direction.
Equality holds when $\xvec$ aligns with the top eigenvector of the
$\kout_{u\to\mathrm{UE}}$-restricted exposure operator.
\end{theorem}

\noindent Sketch.  Apply \texttt{rib.tex} \S\,7 (per-direction
pseudo-Brewster locking) to replace
$\Qbist(\thetab,\kout)\preceq \frac{1-T_0}{T_0}\Qabs^{(\kout)}(\thetab)$.
Bound $\Qabs^{(\kout)}(\thetab)\preceq
\mathcal{V}_{\mathrm{out}}(\thetab)\,\Qabs(\thetab)$ since the
$\kout$-visibility gate restricts integration to a fraction
$\mathcal{V}_{\mathrm{out}}$ of the lit area.  Tightness when $\xvec$
is the top eigenvector follows because both inequalities are tight
on the same eigenvector.  $\square$

\subsection{The gait-action principle}

\Cref{thm:pareto-1} immediately implies a structural statement about
which pose moves are gradient-good:

\begin{corollary}[Gait-action principle]\label{cor:gait}
At a fixed precoder $\xvec$, the pose move that maximises the
body-mediated capacity contribution is the move that simultaneously
\emph{(i)} aligns the phone direction $\kout_{u\to\mathrm{UE}}(\thetab)$
with the top eigenvector of $\Qabs^{(\kout)}(\thetab)$, and
\emph{(ii)} increases $\mathcal{V}_{\mathrm{out}}(\thetab)$ by
exposing more lit area to the phone.  Both effects increase
$P_{\mathrm{ab}}(\thetab,\xvec)$ proportionally, so the body-mediated
capacity-per-dose ratio is invariant.
\end{corollary}

In words: the user can buy capacity but cannot buy a better
capacity-per-dose ratio.  This is what makes the dose receipt honest:
no pose move ``smuggles in'' rate without paying its proportional
price in absorbed power.  The trade is body-universal (driven by
$T_0$), and the scalar $T_0/(1-T_0)\approx 0.32$ is the conversion
rate at \SI{28}{GHz}.

\begin{remark}[Direct path is exempt]\label{rem:los-exempt}
\Cref{thm:pareto-1} bounds only the body-mediated contribution
$P_{\mathrm{body}}$; the LOS contribution to receive power is not
bounded by absorbed dose.  Pose moves that improve the LOS path
(e.g., un-occluding the line of sight to the BS) are gradient-good
without invoking the trade above.  In a body-shadowed scenario this
is the dominant gradient direction, and $P_{\mathrm{body}}$ is a
secondary effect.  The trade matters in the regime where the LOS is
already open and the user is attempting to extract additional
multipath rank from their own body.
\end{remark}

\Cref{rem:los-exempt} matters because it tells us where each gradient
component comes from.  The hero scenario in the companion is a
body-shadowed LOS at a window: the user takes a few-degree torso
turn that opens the LOS, and the gradient is dominated by the
unshadowing term, not the bistatic gain.  This is good for
storytelling: the user's intuition (``turn so my body isn't between
me and the window'') is exactly what the gradient says.

% ============================================================
\section{Closed-loop architecture}\label{sec:loop}
% ============================================================

The mathematics is the same regardless of where you put the gradient
computation.  Two arrangements are worth distinguishing.

\subsection{Where the gradient lives}

\paragraph{Network-side.}  The BS computes
$\nabla_\thetab\mathcal{L}$ using the path dictionary it already has
plus a copy of the user's SMPL-X shape $\betab$ and the user's
current pose estimate $\thetab^{(0)}$.  Suggested pose $\thetab^*$ is
sent to the user.

\paragraph{User-side.}  The phone computes
$\nabla_\thetab\mathcal{L}$ using a path dictionary the BS sends it,
plus the on-device SMPL-X model.  Suggested pose $\thetab^*$ is
displayed locally.

The user-side variant has the privacy advantage that the user's
shape and pose never leave the device, only the resulting move
suggestion or its acceptance.  The network-side variant requires
the user to send $\betab,\thetab^{(0)}$ but spares the device the
\SI{15}{ms} per-slot gradient computation.  We will state the
information-theoretic budget for both.

\subsection{Information channels and rates}

The four channels in the loop, with order-of-magnitude rates:

\begin{center}
\begin{tabular}{lll}
\toprule
\textbf{Channel} & \textbf{Payload} & \textbf{Rate} \\
\midrule
DL: path dictionary               & $\sim 50$ paths $\times$ 4 floats        & \SI{16}{kbps} \\
DL: pose suggestion (network-side) & 72 floats $\times$ \SI{1}{Hz}            & \SI{2.3}{kbps} \\
UL: pose estimate $\thetab^{(0)}$  & 72 floats $\times$ \SI{1}{Hz}            & \SI{2.3}{kbps} \\
UL: shape $\betab$ (one-time)      & 10 floats once per session              & \SI{40}{B} \\
\bottomrule
\end{tabular}
\end{center}

The total bidirectional control bandwidth is
$\le$\SI{20}{kbps}, four orders of magnitude below the data rate
the loop negotiates over.  The Hochwald-style ``do not waste
spectrum on control'' concern is moot here; the loop is essentially
free.

\subsection{Latency budget}

The dominant latency is the user's own actuation: the time between
the BS sending a pose suggestion and the user completing the move,
\SI{200}{ms} to \SI{2}{s}.  This sets the loop period.  Within one
loop cycle, the BS can re-precode at the slot cadence (\SI{0.1}{ms})
and re-evaluate the gradient at the pose-update cadence
(\SI{15}{ms} on a phone-class GPU, \SI{1}{ms} on the BS).  No
hard real-time constraint binds.

\subsection{Privacy by gradient passing}\label{sec:privacy}

In the user-side variant, the phone only ever transmits to the BS:
\emph{(i)} the current pose estimate, \emph{(ii)} the suggestion
acceptance flag, and optionally \emph{(iii)} an aggregated
\emph{absorbed-dose receipt} of the slot.  The body-shape vector
$\betab$ stays on-device; the gradient $\nabla_\thetab\mathcal{L}$
is computed on-device using the BS's path dictionary as the only
network-supplied input.  This is structurally similar to federated
learning: the heavy data (the body's geometry) stays local, and
only summary statistics flow.  The threat surface is the path
dictionary, which is identifying only of the array's view of the
user's location, which the BS already has.

\begin{remark}[Adversarial pose suggestions]
A compromised BS could in principle send pose suggestions that
maximise dose rather than capacity.  The user-side gradient
computation defends against this trivially: the phone re-derives
the suggestion from the loss \eqref{eq:loss} and discards
suggestions that disagree.  The discrepancy bound is the path
dictionary's calibration noise, $\sim$\SI{1}{dB}; suggestions that
exceed this discrepancy are flagged.  This is a deployment-level
detail and not the paper's concern.
\end{remark}

% ============================================================
\section{Convergence and well-posedness}\label{sec:convergence}
% ============================================================

The joint problem is non-convex.  The loss $\mathcal{L}$ is a sum of
a smooth concave-in-$\xvec$ (the rate term) plus a smooth
quasiconvex-in-$\xvec$ exposure penalty plus a strictly convex
quadratic in $\thetab$.  The alternating scheme is a coordinate
descent on a non-convex bivariate.  We claim weak guarantees only,
which is what we have a right to in this regime.

\begin{proposition}[Monotone descent]\label{prop:descent}
Under \cref{thm:Q-diff}, the alternating scheme of \cref{sec:loss}
produces a non-increasing sequence
$\mathcal{L}(\thetab^{(\ell)},\xvec^{(\ell)})$ with limit a
stationary point of the loss.
\end{proposition}

\noindent Sketch.  Each inner step solves the precoder QCQP
exactly (or by warm-started projected ZF in the operational regime;
see JSAC v5).  Each outer step is one gradient step on a Lipschitz
$\nabla_\thetab\mathcal{L}$ with Armijo line search.  Both steps
are descent steps; the sum is bounded below ($-R\ge -R_{\max}$,
$\lambda_E[\cdot]_+\ge 0$, $\lambda_MC\ge 0$); convergence to a
stationary point follows from the standard descent lemma.  $\square$

\begin{remark}[Local versus global]
The objective is non-convex in $\thetab$: a body-shadowed LOS
typically has two local optima corresponding to ``turn left to
unshadow'' and ``turn right to unshadow''.  We do not pursue global
optimisation; the user picks the side they prefer based on their
seating, and the gradient takes them to the nearest local optimum.
A multi-start scheme is straightforward but cosmetic.
\end{remark}

\begin{remark}[Saddle-free]
There are no saddle escapes to worry about because the precoder
inner problem is exactly solvable; we are not in a saddle-point
formulation \emph{per se}, just an alternating descent.  Compare
to GAN-style min-max where the inner problem is also non-convex
and saddle escape is delicate.  Here the inner problem is convex;
the outer non-convexity is over $\thetab$ alone and is benign in
the body-shadowed regime.
\end{remark}

% ============================================================
\section{Connection to the JSAC2 SI scope}\label{sec:si}
% ============================================================

The SI keywords map to this construction in the following one-to-one
way.  We collect them here so the introduction of any submission
can simply quote.

\begin{enumerate}[leftmargin=1.6em,itemsep=3pt]
\item \emph{Closed-loop optimization.}  \Cref{sec:loop} is a
literal closed loop with measurable feedback, computable
gradients, and bounded latency.
\item \emph{Differentiable control.}  \Cref{thm:Q-diff} is the
differentiability statement; \cref{sec:loss} is the control loss;
\eqref{eq:Q-jac} is the gradient.
\item \emph{Twin-in-the-loop architectures.}  The body twin is in
the loop in the strong sense: removing it makes the gradient
ill-defined.  Compare to JSAC v5 where removing the twin would
just default to a worst-case scaling.
\item \emph{Cross-layer optimization.}  The objective
\eqref{eq:loss} couples PHY (rate, exposure) with the
biomechanical layer (pose, comfort).  ``Cross-layer'' is a
generous reading of two layers; we will own it.
\item \emph{Application-aware design.}  The application is the
user's data session.  ``Application'' here is not the OSI-stack
application layer but the SI's broader meaning of the user-facing
purpose; the construction is parameterised by what the user is
trying to do.
\item \emph{Predictive DT analytics.}  The phone displays the
current dose receipt and the predicted dose-and-rate trajectory
under the suggested move.  This is predictive analytics on the
DT.
\item \emph{Trustworthy data collection.}  In the user-side
variant of \cref{sec:privacy}, body-shape and pose data never
leave the device.  This is at minimum a stronger privacy
posture than the BS-side collection.
\item \emph{XR-adjacent application.}  Pose-aware UE behaviour is
a primitive in XR systems; the same loop applies to a head-mounted
device whose pose tracking is already exposed by the OS.
\end{enumerate}

% ============================================================
\section{What this is not}\label{sec:notthis}
% ============================================================

A note on scoping.

\begin{enumerate}[leftmargin=1.6em,itemsep=3pt]
\item This is not a JCAS / ISAC paper.  We do not estimate pose from
backscatter; we use the on-device pose estimator (IMU + on-screen
AR + the v5 virtual-IMU pipeline).  The bistatic operator
$\Frib$ appears in \cref{thm:pareto-1} only as a forward predictor
of body-mediated channel, not as a sensing channel.  A follow-up
paper that closes the loop with backscatter pose estimation is
plausible but is its own object.
\item This is not a beamforming paper.  The precoder is the
JSAC v5 projected-ZF closed form; we do not propose a new
precoder.  The novelty is in the outer loop.
\item This is not a pose-estimation paper.  The pose estimator is
borrowed from Madgwick / Mahony / SMPL-X-from-IMU literature.
Our use of pose is entirely downstream of estimation.
\item This is not an RIS paper.  The body is not actively tunable;
its gait is.  The analogy to RIS is informal and is used only to
locate the construction in a familiar literature.  A literal RIS
deployment has tunable phase shifters at sub-millimetre granularity;
a body has none.
\item This is not a regulatory-compliance paper.  The exposure
constraint is the engineering substrate that makes the trade
honest.  We do not certify any deployment to ICNIRP-2020; we
borrow the limits as the reference scaling.
\end{enumerate}

% ============================================================
\section{Loose threads and conjectures}\label{sec:threads}
% ============================================================

In the same format as \texttt{rib.tex} \S\,12 and
\texttt{q\_complement.tex} \S\,9: questions, not claims.

\begin{enumerate}[leftmargin=1.6em,itemsep=3pt]
\item \emph{Multi-user RIHB.}  Two users in the same scene, both
opted in.  Their pose-gradient updates couple through their shared
channel; the natural object is the joint loss
$\mathcal{L}(\thetab_1,\thetab_2,\xvec)$ with one Pareto bound per
user.  Conjecture: the joint Pareto frontier is non-improving
in either user when both follow gradients, but the
per-user-conditioned frontiers can improve, giving a Stackelberg-style
lift.  Worth a small numerical experiment.
\item \emph{Adversarial poses.}  Worst-case pose for a given
precoder is the global maximum of $P_{\mathrm{ab}}(\thetab)$ over a
comfort-bounded ball.  This is a robustness question for the BS:
how much rate must it leave on the table to be safe against pose
the user cannot avoid?  Probably a half-day experiment.
\item \emph{Federated $\betab$.}  An operator that wants
population-level capacity statistics could in principle estimate
$\bar\betab$ across users without seeing any individual $\betab$,
via federated averaging.  Whether the resulting average shape is
useful or just a smoothed mannequin is an open question.
\item \emph{Pose noise as channel noise.}  The pose estimator's
\SI{0.5}{\degree} RMS gives a $\Qabs$-evaluation noise of order
$\sim 1\%$ on $\tr\Qabs$ at static bodies.  In motion this is much
larger.  The natural object is a pose-noise-aware Pareto bound;
the JSAC v5 IMU-vs-T-pose ablation gives the empirical numbers we
need to instantiate it.
\item \emph{Sub-\SI{6}{GHz}.}  The pseudo-Brewster collapse
\texttt{q\_complement.tex} \S\,4 holds at mmWave; below
\SI{6}{GHz} the body's $T_0$ is polarisation-dependent and the
$T_0/(1-T_0)$ slope is no longer body-universal.  The construction
extends but the Pareto becomes per-polarisation.  Whether 5G FR1
is a useful regime for RIHB is unclear; the body-blockage
motivation evaporates because FR1 diffraction routes around
torsos cheaply.
\item \emph{Pose suggestion as channel coding.}  The BS's pose
suggestion is a message in $\Reals^{72}$; the user's response is
a binary accept/reject plus the new pose.  This is a low-rate,
high-latency, error-prone channel.  Treating it as a coding
problem is academic but has a charm: if you let the BS suggest a
\emph{distribution} over poses with priors from the comfort
metric, you recover the rate-distortion view of human-in-the-loop
control.
\item \emph{Pose-conditional Pareto frontier in dB.}  The JSAC v5
$\alpha$-vs-WC gap was an aspirational quantitative result that
never landed.  The pose-conditional version is concrete and
measurable: for each pose in the ball $C\le\rho$, plot
$(\log R(\thetab,\xvec^*),P_{\mathrm{ab}}(\thetab,\xvec^*))$ and
compute the dB gap to the worst-case envelope at the same comfort
budget.  This is the headline figure of the companion experiment.
\item \emph{Comfort metric calibration.}  The
weights $c_i$ in \eqref{eq:movement-cost} are the only thing in
this construction that is not derived from physics.  Calibrating
them against a small user study (5 users, ergonomic score on a
dozen pose moves) is the kind of mixed-methods step a careful
referee will ask for, and is a reasonable companion-paper item.
\item \emph{Phone in the body's near field.}  The phone is at
$\sim$\SI{30}{cm} from the wrist, and the wrist is on the body.
The body is in the BS's far field at \SI{30}{m} away
($D^2/\lambda\sim\SI{300}{m}$ but the relevant near-field is the
body's, not the BS's).  The phone-to-body geometry needs the
near-field correction of \texttt{rib.tex} \S\,12 (Fresnel-Kirchhoff
kernel); this is in $\Frib$ and not yet plumbed through.  A real
fix, but mechanical.
\item \emph{The negotiation as a game.}  An operator that bills
on data delivered while a user that pays in pose deviation is a
two-player game with a known utility for each.  The Nash
equilibrium of this game is, conjecturally, the operating point
the closed loop converges to.  Pricing the pose move is
business-model territory and is left to the negotiation.
\end{enumerate}

% ============================================================
\section{The companion experiment}\label{sec:comp}
% ============================================================

The empirical content of the spine is a single hero figure: in a
body-shadowed line-of-sight scenario, the dB-domain capacity-vs-pose
sweep around a baseline $\thetab^{(0)}$, with iso-exposure contours
on the same axes.  See \texttt{JSAC2/scenarios.md} for the
scenario specification, \texttt{JSAC2/poc/} for the proof-of-concept
script that generates the figure, and \texttt{JSAC2/critique\_v1.md}
for the self-audit that motivates the next iteration of the spine.

The headline number we want from the figure: the dB-difference
between baseline rate and gradient-optimal rate at a comfort budget
of \SI{5}{\degree} torso rotation.  My pre-experiment estimate is
\SI{12}{dB} to \SI{18}{dB}.  Anything under \SI{6}{dB} kills the
spine.  Anything over \SI{12}{dB} gives the paper a teeth.

% ============================================================
\section{Summary}\label{sec:sum}
% ============================================================

In one sentence: \emph{the body twin in the JSAC paper is a noun;
this note makes it a verb.}  The same surface-integral operators
that already estimate dose now flow gradients to the user, who is
free to take or ignore them, and whose body becomes the channel-shaping
degree of freedom the network was missing.  The differentiability
contract is small (\cref{thm:Q-diff}); the Pareto bound that makes
the trade honest is body-universal (\cref{thm:pareto-1}); the
information-theoretic budget of the loop is negligible
(\cref{sec:loop}).  What remains is empirical: how much capacity
does a few-degree pose move buy, and at what dose, in the body-shadowed
LOS regime.  The companion experiment in \cref{sec:comp} answers
that question; the spine of any JSAC2 SI submission rests on the
answer.

\end{document}
